% Canonical TeX edition. Initial text inherited from 3b067cc2eba4ac2f2f91248448727639d8ac7c5a
% Source: research/case_study/phase277_coprime_obstruction.md
% Subsequent mathematical edits require a new audit.
% Options for packages loaded elsewhere
\PassOptionsToPackage{unicode}{hyperref}
\PassOptionsToPackage{hyphens}{url}
\documentclass[
  11pt,
  a4paper,
]{article}
\usepackage{xcolor}
\usepackage[margin=1in]{geometry}
\usepackage{amsmath,amssymb}
\setcounter{secnumdepth}{-\maxdimen} % remove section numbering
\usepackage{iftex}
\ifPDFTeX
  \usepackage[T1]{fontenc}
  \usepackage[utf8]{inputenc}
  \usepackage{textcomp} % provide euro and other symbols
\else % if luatex or xetex
  \usepackage{unicode-math} % this also loads fontspec
  \defaultfontfeatures{Scale=MatchLowercase}
  \defaultfontfeatures[\rmfamily]{Ligatures=TeX,Scale=1}
\fi
\usepackage{lmodern}
\ifPDFTeX\else
  % xetex/luatex font selection
  \setmainfont[]{STIXGeneral}
  \setmonofont[]{Menlo}
\fi
% Use upquote if available, for straight quotes in verbatim environments
\IfFileExists{upquote.sty}{\usepackage{upquote}}{}
\IfFileExists{microtype.sty}{% use microtype if available
  \usepackage[]{microtype}
  \UseMicrotypeSet[protrusion]{basicmath} % disable protrusion for tt fonts
}{}
\makeatletter
\@ifundefined{KOMAClassName}{% if non-KOMA class
  \IfFileExists{parskip.sty}{%
    \usepackage{parskip}
  }{% else
    \setlength{\parindent}{0pt}
    \setlength{\parskip}{6pt plus 2pt minus 1pt}}
}{% if KOMA class
  \KOMAoptions{parskip=half}}
\makeatother
\setlength{\emergencystretch}{3em} % prevent overfull lines
\providecommand{\tightlist}{%
  \setlength{\itemsep}{0pt}\setlength{\parskip}{0pt}}
\usepackage{newunicodechar}
\newfontfamily\numeralfallback{STIX Two Text}
\newunicodechar{∎}{\ensuremath{\blacksquare}}
\usepackage{fvextra}
\DefineVerbatimEnvironment{Highlighting}{Verbatim}{breaklines,breakanywhere,breakanywheresymbolpre={},breakanywheresymbolpost={},breaksymbolleft={},breaksymbolright={},commandchars=\\\{\}}
\newunicodechar{⁰}{{\numeralfallback\char"2070\relax}}
\newunicodechar{₀}{{\numeralfallback\char"2080\relax}}
\newunicodechar{₁}{{\numeralfallback\char"2081\relax}}
\newunicodechar{₂}{{\numeralfallback\char"2082\relax}}
\newunicodechar{₃}{{\numeralfallback\char"2083\relax}}
\newunicodechar{₄}{{\numeralfallback\char"2084\relax}}

\usepackage{seqsplit}
\renewcommand{\texttt}[1]{{\ttfamily\seqsplit{#1}}}
\setlength{\emergencystretch}{6em}

\newunicodechar{⁻}{{\numeralfallback\char"207B\relax}}
\newunicodechar{⁸}{{\numeralfallback\char"2078\relax}}
\newunicodechar{ⁱ}{{\numeralfallback\char"2071\relax}}
\newunicodechar{⁴}{{\numeralfallback\char"2074\relax}}
\newunicodechar{⁷}{{\numeralfallback\char"2077\relax}}
\newunicodechar{⁶}{{\numeralfallback\char"2076\relax}}
\newunicodechar{⁵}{{\numeralfallback\char"2075\relax}}
\newunicodechar{₈}{{\numeralfallback\char"2088\relax}}
\newunicodechar{ᵢ}{\textsubscript{i}}
\newunicodechar{ⱼ}{\textsubscript{j}}
\newunicodechar{ₜ}{\textsubscript{t}}
\usepackage{bookmark}
\IfFileExists{xurl.sty}{\usepackage{xurl}}{} % add URL line breaks if available
\urlstyle{same}
% fallback for those not using the hyperref driver hyperxmp:
\makeatletter
\@ifundefined{xmpquote}{\newcommand{\xmpquote}[1]{#1}}{}
\makeatother
\hypersetup{
  hidelinks,
  pdfcreator={LaTeX via pandoc}}

\author{}
\date{}

\begin{document}

\section{A necessary condition for a coprime
kernel}\label{a-necessary-condition-for-a-coprime-kernel}

Private continuation, phase 277. This obstruction concerns independent
marked groups in arbitrary rank. The displayed example is expressly not
a polytope monodromy group, and is not a counterexample to the
finite-cover conjecture for polytopes.

Let \(\pi:H\to G\) be an epimorphism of finite marked string groups, and
suppose \(H\) has the intersection property. Take disjoint color sets
\(I,J\) such that every mark in one set commutes with every mark in the
other; it is enough that every cross distance exceeds one. Put
\(A=H_I\), \(B=H_J\), and \(K=\ker\pi\).

\textbf{Proposition 1.} The order of \(G_I\cap G_J\) divides \(|K|\).
Consequently a C-cover with kernel order coprime to \(|G|\) is possible
only if all such separated parabolic intersections in \(G\) are trivial.

\textbf{Proof.} The subgroups \(A\) and \(B\) commute and intersect
trivially, so multiplication identifies \(AB\) with \(A\times B\). Let
\(D=\pi(A)\cap\pi(B)\). Define a homomorphism \[
 f:K\cap AB\longrightarrow D,\qquad f(ab)=\pi(a).
\] The factors \(a,b\) are unique. If \(ab\in K\), then
\(\pi(a)=\pi(b)^{-1}\in D\), so the map is well-defined. Commutation of
\(A\) and \(B\) proves multiplicativity. For every \(d\in D\), choose
\(a\in A\) and \(b\in B\) with \(\pi(a)=d=\pi(b)\). Then
\(ab^{-1}\in K\cap AB\) and \(f(ab^{-1})=d\), proving surjectivity. The
first isomorphism theorem and Lagrange's theorem now show
\(|D|\mid|K\cap AB|\mid|K|\). Since \(|D|\) also divides \(|G|\), the
coprime assertion follows. ∎

This condition cannot be omitted from a higher-rank extension of the
rank-four p-kernel theorem. Let
\(D_8=\langle a,b\mid a^2=b^2=(ab)^4=1\rangle\), with central involution
\(z=(ab)^2\). Use two copies, with marks \((a_0,a_1)\) and
\((b_3,b_4)\), and define \[
 G=\bigl(D_8\times D_8\bigr)/\langle(z,z)\rangle\times C_2.
\] Mark the five factors in the order \(a_0,a_1,t,b_3,b_4\), where \(t\)
generates the final direct factor. Then \(|G|=64\). The five marks are
independent: the abelianization of the central product is \(C_2^4\),
since the identified central subgroup already lies in the commutator
subgroup, and the extra factor contributes one more independent
coordinate. All required string commutations hold.

Both maximal contiguous end groups have order 32 and are isomorphic,
with the indicated separated components, to \(D_8\times C_2\times C_2\).
They are C-groups; so are all their standard parabolics. But \[
 G_{\{0,1\}}\cap G_{\{3,4\}}=\langle\bar z\rangle\cong C_2.
\] Proposition 1 implies that \textbf{every finite C-cover of this
marked group has even kernel order}. In particular no odd-prime p-group
kernel can work, regardless of its dimension or number of stages. The
product \(D_8\times D_8\times C_2\) itself is a finite C-cover with
kernel of order2, so the example does not obstruct finite existence.

Why this is not a polytope counterexample is equally precise. In the
monodromy action of an actual polytope, a word using only colors in
\(I\) fixes every face of ranks outside \(I\), at every flag. If an
element also belongs to the subgroup on a disjoint set \(J\), it fixes
faces of every rank, and therefore fixes every flag. The monodromy
permutation is the identity. Thus actual monodromy satisfies \[
 G_I\cap G_J=1\quad\text{whenever }I\cap J=\varnothing,
\] with no need for the separated-color assumption. The example above
violates this necessary condition.

The remaining finite-polytope problem must exploit this genuine
flag-action condition when attempting a uniform p-kernel construction in
higher rank. Extreme independence alone, sufficient in rank four, no
longer suffices. No assertion of historical novelty for Proposition 1 or
this central-product example is made without a separate literature
check.

\subsection{Exact finite model and a qualification about
covers}\label{exact-finite-model-and-a-qualification-about-covers}

The independent reader can reconstruct the example without identifying a
small-group catalogue entry. Write an element of the order-64 group as
\((v_0,v_1,v_3,v_4,z,t)\in\mathbb F_2^6\), with multiplication \[
 (v,z,t)(w,d,u)=
 (v+w,\ z+d+v_1w_0+v_4w_3,\ t+u).
\] The five marks are the coordinate vectors in positions
\(v_0,v_1,t,v_3,v_4\). The order-128 cover uses two central coordinates,
with cocycles \(v_1w_0\) and \(v_4w_3\) separately; its quotient map
adds those coordinates. The program
\texttt{phase277\_obstruction\_check.py} checks both multiplication
tables, all associativity triples, the marked epimorphism and its
two-element kernel, all 1,024 parabolic intersections of the cover, and
all 256 pairs in each end of the base group. This is a certificate for
the displayed finite model. The divisibility assertion about every
finite C-cover is Proposition 1's written proof, not a conclusion of
those finite checks.

There is no claim that the order-128 cover is the only minimal cover. In
particular, centrality of \(g_2\) does not force centrality of its lift:
the adjacent commutations with \(g_1\) and \(g_3\) need not survive in a
cover. It would be incorrect to split every cover into three direct
factors on the basis of the direct-product presentation of this
particular base.

\end{document}
